\section{Zarządzanie przedsięwzięciami informatycznymi}

\subsection{Czym charakteryzuje się projekt informatyczny?}

Projekt informatyczny jest tymczasowym przedsięwzięciem prowadzącym do wytworzenia unikalnego produktu, usługi albo rezultatu związanego z systemem informatycznym. Może to być stworzenie aplikacji, wdrożenie systemu ERP, migracja danych, integracja usług, modernizacja infrastruktury, budowa hurtowni danych albo wdrożenie modelu uczenia maszynowego.

\subsubsection{Cechy projektu}

Projekt ma:
\begin{itemize}
  \item jasno określony cel biznesowy albo techniczny,
  \item początek i koniec,
  \item ograniczony budżet, czas i zakres,
  \item interesariuszy,
  \item zespół projektowy,
  \item ryzyka i niepewność,
  \item plan realizacji,
  \item kryteria akceptacji.
\end{itemize}

W IT szczególnie istotne są niepewność wymagań, zmienność technologii, integracje z istniejącymi systemami, zależność od kompetencji ludzi oraz trudność szacowania prac twórczych.

\subsubsection{Potrójne ograniczenie i jego rozszerzenia}

Klasycznie projekt opisuje się przez trójkąt:
\[
\text{zakres} \quad-\quad \text{czas} \quad-\quad \text{koszt}.
\]
Zmiana jednego elementu wpływa na pozostałe. Jeśli zakres rośnie przy stałym terminie, trzeba zwiększyć zasoby, obniżyć jakość albo zaakceptować ryzyko opóźnienia.

W praktyce dodaje się jakość, ryzyko, satysfakcję klienta, wartość biznesową i ograniczenia organizacyjne. Projekt IT nie jest sukcesem tylko dlatego, że kod działa. Sukces oznacza dostarczenie wartości w akceptowalnym koszcie, czasie i jakości.

\subsubsection{Specyfika projektów informatycznych}

Projekty informatyczne mają kilka cech odróżniających je od wielu projektów inżynierii materialnej:
\begin{description}
  \item[Niematerialny produkt] Oprogramowania nie widać bezpośrednio, dlatego postęp bywa trudny do oceny, jeśli nie ma działających przyrostów.
  \item[Zmienność wymagań] Użytkownicy często doprecyzowują potrzeby dopiero po zobaczeniu prototypu.
  \item[Wysoka złożoność zależności] Integracje, dane, infrastruktura, bezpieczeństwo i wydajność tworzą zależności ukryte.
  \item[Niski koszt kopiowania, wysoki koszt wytworzenia] Samo powielenie programu jest tanie, ale jego zaprojektowanie, testowanie i utrzymanie jest kosztowne.
  \item[Dług techniczny] Decyzje przyspieszające krótkoterminowo mogą zwiększać koszt dalszego rozwoju.
  \item[Trudność szacowania] Zadania programistyczne są często poznawczo nowe, więc estymacje mają dużą niepewność.
\end{description}

\subsubsection{Interesariusze}

Interesariusz to osoba albo organizacja wpływająca na projekt albo dotknięta jego wynikiem. W projekcie IT są to zwykle:
\begin{itemize}
  \item sponsor biznesowy,
  \item product owner albo właściciel produktu,
  \item użytkownicy końcowi,
  \item zespół projektowy,
  \item administratorzy i zespoły utrzymaniowe,
  \item dział bezpieczeństwa,
  \item compliance i prawnicy,
  \item dostawcy zewnętrzni,
  \item klienci.
\end{itemize}

Dobre zarządzanie projektem wymaga rozumienia, kto podejmuje decyzje, kto akceptuje rezultat, kto dostarcza wymagania i kto ponosi skutki wdrożenia.

\subsubsection{Zakres, wymagania i akceptacja}

Zakres określa, co projekt ma dostarczyć. Wymagania mogą być:
\begin{itemize}
  \item funkcjonalne - co system ma robić,
  \item niefunkcjonalne - wydajność, bezpieczeństwo, dostępność, użyteczność,
  \item biznesowe - jaki problem organizacji ma zostać rozwiązany,
  \item techniczne - standardy, technologie, integracje,
  \item regulacyjne - zgodność z prawem i politykami.
\end{itemize}

Kryteria akceptacji muszą być weryfikowalne. Zamiast "system ma być szybki" lepiej zapisać "95 percent odpowiedzi API dla operacji X ma czas poniżej 300 ms przy obciążeniu Y".

\subsubsection{Ryzyka typowe dla projektów IT}

\begin{itemize}
  \item niedoszacowanie złożoności,
  \item niestabilne wymagania,
  \item brak dostępności kluczowych osób,
  \item opóźnienia integracji,
  \item niska jakość danych,
  \item problemy z bezpieczeństwem,
  \item vendor lock-in,
  \item brak środowisk testowych,
  \item niedostateczne testy,
  \item brak planu migracji i rollbacku,
  \item konflikt między terminem a jakością.
\end{itemize}

\subsubsection{Co powiedzieć na obronie}

Projekt informatyczny to tymczasowe przedsięwzięcie tworzące unikalny rezultat IT. Charakteryzuje się określonym celem, zakresem, terminem, budżetem, interesariuszami i ryzykiem. Szczególne dla IT są zmienne wymagania, niematerialność produktu, trudność szacowania, złożoność integracji i dług techniczny.

\subsection{Główne różnice między projektem a pracą operacyjną}

Projekt i praca operacyjna są potrzebne organizacji, ale mają różny charakter. Projekt wprowadza zmianę; operacje utrzymują ciągłe działanie.

\subsubsection{Porównanie}

\begin{longtable}{p{0.25\textwidth}p{0.33\textwidth}p{0.33\textwidth}}
\toprule
Kryterium & Projekt & Praca operacyjna \\
\midrule
Czas trwania & tymczasowy, ma początek i koniec & ciągła, powtarzalna \\
Cel & wytworzenie unikalnego rezultatu albo zmiany & utrzymanie działania i świadczenie usług \\
Zakres & zdefiniowany dla konkretnego przedsięwzięcia & stabilny, procesowy \\
Niepewność & zwykle wysoka & zwykle niższa, procedury są znane \\
Sukces & dostarczenie wartości w czasie, budżecie i jakości & stabilność, efektywność, SLA, jakość obsługi \\
Zespół & często czasowy i interdyscyplinarny & stały albo rotacyjny \\
Przykład IT & budowa nowej aplikacji & utrzymanie aplikacji produkcyjnej \\
\bottomrule
\end{longtable}

\subsubsection{Relacja między projektem i operacjami}

Projekt często przekazuje wynik do operacji. Przykład: zespół projektowy wdraża system, a potem zespół utrzymaniowy monitoruje go, obsługuje incydenty, wykonuje aktualizacje i wspiera użytkowników.

Granica nie zawsze jest ostra. Rozwój produktu cyfrowego może być ciągłym strumieniem zmian, ale poszczególne większe inicjatywy nadal mają charakter projektowy. W podejściu produktowym mówi się częściej o ciągłym rozwoju produktu niż o jednorazowych projektach, jednak nadal występują ograniczone inicjatywy, budżety i cele.

\subsubsection{Przykład}

Wdrożenie modułu płatności online jest projektem: ma zakres, termin, integracje, testy bezpieczeństwa i odbiór. Codzienne monitorowanie transakcji, reagowanie na błędy płatności i aktualizacja certyfikatów to operacje.

\subsubsection{Co powiedzieć na obronie}

Projekt jest tymczasowy i tworzy unikalną zmianę, a praca operacyjna jest ciągła i powtarzalna. Projekt kończy się przekazaniem rezultatu, natomiast operacje utrzymują usługę, proces albo system. W IT projekty często budują lub zmieniają systemy, a operacje zapewniają ich stabilne działanie.

\subsection{Metodyki zarządzania projektami: charakterystyka, różnice i cechy wspólne}

Metodyka zarządzania projektem określa sposób planowania, organizacji, realizacji, kontroli i zamykania prac. W IT najczęściej porównuje się podejścia predykcyjne, zwinne i hybrydowe.

\subsubsection{Waterfall}

Waterfall, czyli model kaskadowy, zakłada sekwencyjne fazy:
\begin{enumerate}
  \item analiza wymagań,
  \item projekt,
  \item implementacja,
  \item testowanie,
  \item wdrożenie,
  \item utrzymanie.
\end{enumerate}

Zalety:
\begin{itemize}
  \item jasny plan i dokumentacja,
  \item dobre dopasowanie do stałych wymagań,
  \item łatwiejsze kontraktowanie zakresu,
  \item formalna kontrola zmian.
\end{itemize}

Wady:
\begin{itemize}
  \item późna informacja zwrotna od użytkownika,
  \item kosztowne zmiany wymagań,
  \item ryzyko odkrycia problemów dopiero pod koniec,
  \item słabe dopasowanie do niepewnych produktów cyfrowych.
\end{itemize}

Waterfall jest sensowny, gdy wymagania są stabilne, środowisko regulowane, a koszt zmian wysoki, np. wybrane systemy krytyczne albo kontrakty publiczne.

\subsubsection{PMBOK i podejście procesowe}

PMBOK nie jest jedną sekwencją prac, lecz zbiorem dobrych praktyk zarządzania projektami. Obejmuje obszary takie jak zakres, harmonogram, koszt, jakość, zasoby, komunikacja, ryzyko, zamówienia i interesariusze. Jest metodycznie neutralny: można go stosować w projektach kaskadowych, zwinnych i hybrydowych.

Mocną stroną PMBOK jest kompletność zarządcza. Słabością w małych zespołach może być nadmiar formalizacji, jeśli stosuje się go bez dostosowania.

\subsubsection{PRINCE2}

PRINCE2 jest metodyką silnie nastawioną na uzasadnienie biznesowe, role, produkty zarządcze i kontrolę etapową. Kluczowe elementy:
\begin{itemize}
  \item ciągła zasadność biznesowa,
  \item uczenie się z doświadczeń,
  \item zdefiniowane role i odpowiedzialności,
  \item zarządzanie etapami,
  \item zarządzanie przez wyjątki,
  \item koncentracja na produktach,
  \item dostosowanie do środowiska projektu.
\end{itemize}

PRINCE2 dobrze pasuje do organizacji wymagających kontroli i raportowania, ale może być zbyt formalny dla małych, szybko zmieniających się produktów.

\subsubsection{Agile}

Agile to rodzina podejść opartych na iteracyjności, przyrostowym dostarczaniu wartości, współpracy z klientem i reagowaniu na zmianę. Nie oznacza braku planu. Oznacza planowanie adaptacyjne i częstą weryfikację założeń.

Wspólne idee:
\begin{itemize}
  \item krótkie iteracje,
  \item częste dostarczanie działającego przyrostu,
  \item bliska współpraca z użytkownikiem lub właścicielem produktu,
  \item priorytetyzowany backlog,
  \item retrospekcje i ciągłe usprawnianie,
  \item akceptacja zmienności wymagań.
\end{itemize}

\subsubsection{Scrum}

Scrum jest ramą pracy zwinnej. Role:
\begin{description}
  \item[Product Owner] odpowiada za wartość produktu i backlog.
  \item[Scrum Master] dba o proces, usuwa przeszkody i wspiera zespół.
  \item[Developers] wytwarzają przyrost produktu.
\end{description}

Zdarzenia:
\begin{itemize}
  \item Sprint,
  \item Sprint Planning,
  \item Daily Scrum,
  \item Sprint Review,
  \item Sprint Retrospective.
\end{itemize}

Artefakty:
\begin{itemize}
  \item Product Backlog,
  \item Sprint Backlog,
  \item Increment.
\end{itemize}

Scrum dobrze działa, gdy produkt można rozwijać przyrostowo, a wymagania są zmienne. Nie rozwiązuje sam problemów technicznych: wymaga praktyk inżynierskich, testów, CI/CD i dbania o architekturę.

\subsubsection{Kanban}

Kanban koncentruje się na przepływie pracy. Typowe zasady:
\begin{itemize}
  \item wizualizacja pracy na tablicy,
  \item limity WIP, czyli ograniczenie pracy w toku,
  \item zarządzanie przepływem,
  \item jawne polityki procesu,
  \item ciągłe doskonalenie.
\end{itemize}

Kanban jest dobry dla utrzymania, zespołów operacyjnych, obsługi zgłoszeń i środowisk, gdzie praca napływa nieregularnie. W przeciwieństwie do Scruma nie wymaga sprintów.

\subsubsection{Extreme Programming}

XP jest metodyką zwinną silnie skupioną na praktykach inżynierskich:
\begin{itemize}
  \item pair programming,
  \item test-driven development,
  \item continuous integration,
  \item refaktoryzacja,
  \item prosta architektura,
  \item małe wydania,
  \item wspólna własność kodu.
\end{itemize}

XP odpowiada na lukę często widoczną w Scrumie: sama organizacja iteracji nie gwarantuje jakości technicznej.

\subsubsection{Lean}

Lean skupia się na maksymalizacji wartości i eliminacji marnotrawstwa. W IT marnotrawstwem może być nadmiar funkcji, oczekiwanie, przełączanie kontekstu, częściowo wykonana praca, defekty, ręczne czynności i nadmierna biurokracja. Lean podkreśla optymalizację całego strumienia wartości, nie tylko lokalną produktywność osób.

\subsubsection{DevOps}

DevOps nie jest klasyczną metodyką projektową, ale zestawem praktyk łączących rozwój, utrzymanie i automatyzację dostarczania. Ważne elementy:
\begin{itemize}
  \item CI/CD,
  \item infrastructure as code,
  \item monitoring i observability,
  \item automatyzacja testów,
  \item szybki rollback,
  \item współodpowiedzialność za produkcję.
\end{itemize}

DevOps skraca dystans między projektem i operacjami.

\subsubsection{Podejścia hybrydowe}

W praktyce często łączy się elementy. Przykład: formalne etapy i budżet według PRINCE2, a realizacja zespołu w Scrumie. Albo roadmapa produktu i kwartalne cele, ale codzienny przepływ w Kanbanie. Hybryda jest sensowna, jeśli wynika z potrzeb organizacji, a nie z niekonsekwencji.

\subsubsection{Cechy wspólne metodyk}

Mimo różnic większość metodyk obejmuje:
\begin{itemize}
  \item definiowanie celu i zakresu,
  \item role i odpowiedzialności,
  \item planowanie pracy,
  \item monitorowanie postępu,
  \item zarządzanie ryzykiem,
  \item komunikację z interesariuszami,
  \item kontrolę jakości,
  \item obsługę zmian,
  \item zamknięcie albo przekazanie rezultatu.
\end{itemize}

\subsubsection{Najważniejsze różnice}

Waterfall zakłada stabilność wymagań i sekwencyjność. Agile zakłada zmienność i iteracyjne dostarczanie wartości. Scrum daje ramę pracy zespołowej w iteracjach. Kanban optymalizuje przepływ i limity pracy w toku. PRINCE2 i PMBOK są bardziej zarządcze i formalne. XP skupia się na jakości inżynierskiej. DevOps skupia się na automatyzacji dostarczania i odpowiedzialności za działanie systemu.

\subsubsection{Co powiedzieć na obronie}

Metodyki różnią się podejściem do zmiany, planowania i kontroli. Kaskadowe są dobre przy stabilnych wymaganiach, zwinne przy niepewności i potrzebie częstej informacji zwrotnej, Kanban przy ciągłym przepływie zadań, PRINCE2/PMBOK przy formalnym zarządzaniu, XP przy nacisku na praktyki programistyczne. Wspólne są cel, role, planowanie, monitorowanie, ryzyko, jakość i komunikacja.
